% This template is used for all validation runs in the report (which can use different input datasets!)
\subsection{Validation run \detokenize{$<ContentVars['report_run_index']>$}}


\subsubsection{Validation Settings}

The following validation settings were used in the comparison against ERA5-Land.
For a full list of used settings, see the validation run results page for this run via the provided link in Table \ref{tab:val_run_list}.

\begin{itemize}
    \item Spatial reference (spatial matching): \textbf{\detokenize{$<ContentVars['ConfigVars']['spatial_reference']>$}}
    \item Temporal reference (temporal matching): \textbf{\detokenize{$<ContentVars['ConfigVars']['temporal_reference']>$}}
    \item Temporal matching window was set to \textbf{\detokenize{$<ContentVars['ConfigVars']['temporal_matching']>$} hours}.
    \item Scaling reference: \textbf{\detokenize{$<ContentVars['ConfigVars']['scaling_reference']>$}}
    \item Validation metrics \detokenize{$<'\\textbf{are}' if ContentVars['ConfigVars']['anomalies_method']!='none' else '\\textbf{are not}'>$} based on anomalies.
    \item Triple collocation analysis \detokenize{$<'\\textbf{was}' if ContentVars['ConfigVars']['metrics'][0]['value'] else '\\textbf{was not}'>$} performed.
\end{itemize}

\subsubsection{Results}

\noindent \ac{SMOS} \ac{L2} was compared against ERA5-Land at \detokenize{$<ContentVars['NetcdfVars']['n_gpis']>$} grid points.
The datasets are spatially matched via a nearest neighbour search, with \detokenize{$<ContentVars['ConfigVars']['spatial_reference']>$} as the spatial reference.
Validation results are therefore computed on the \detokenize{$<ContentVars['ConfigVars']['spatial_reference']>$} grid.

\vspace{5mm}

\noindent Table \ref{tab:sumstats1} lists the summary statistics in terms of number of available observations for the metric computation ("\# Obs."), i.e., matching measurements over time between the satellite and in situ site, the correlation coefficient ("Pearson's R"), Bias (difference in means) and \ac{ubRMSD}.
Additional metrics can be found in the validation results \href{\detokenize{$<ContentVars['NetcdfMetaVars']['url']>$}}{online}.

\begin{table}[htb!]
    \centering
    \caption{Summary statistics of \ac{SMOS} \ac{L2} against ERA5-Land.}
    \label{tab:sumstats2}
    \begin{tabular}{ l r r r r }
        \toprule
        \textbf{Metric} & \textbf{\#Obs.} & \textbf{Pearson's R} & \textbf{Bias}  & \textbf{ubRMSD}  \\
        \midrule
        \textbf{Mean}   & \detokenize{$<ContentVars['SummaryVars']['MEAN_N. OBS.']>$}   & \detokenize{$<ContentVars['SummaryVars']['MEAN_PEARSON'S R']>$}     & \detokenize{$<ContentVars['SummaryVars']['MEAN_BIAS']>$}   & \detokenize{$<ContentVars['SummaryVars']['MEAN_UNBIASED ROOT-MEAN-SQUARE DEVIATION']>$}   \\
        \textbf{Median}   & \detokenize{$<ContentVars['SummaryVars']['MEDIAN_N. OBS.']>$}   & \detokenize{$<ContentVars['SummaryVars']['MEDIAN_PEARSON'S R']>$}     & \detokenize{$<ContentVars['SummaryVars']['MEDIAN_BIAS']>$}   & \detokenize{$<ContentVars['SummaryVars']['MEDIAN_UNBIASED ROOT-MEAN-SQUARE DEVIATION']>$}   \\
        \textbf{IQ-Range}   & \detokenize{$<ContentVars['SummaryVars']['IQ RANGE_N. OBS.']>$}   & \detokenize{$<ContentVars['SummaryVars']['IQ RANGE_PEARSON'S R']>$}     & \detokenize{$<ContentVars['SummaryVars']['IQ RANGE_BIAS']>$}   & \detokenize{$<ContentVars['SummaryVars']['IQ RANGE_UNBIASED ROOT-MEAN-SQUARE DEVIATION']>$}   \\
        \bottomrule
    \end{tabular}
\end{table}

% Here we do some on-the-fly computation with numpy
\noindent For \detokenize{$<np.mean(ContentVars['NetcdfVars']['percent_ok']):.2f>$} \% of all attempted grid cells, validation metrics could be computed.
For the remaining grid cells, an error occurred (compare Fig. \ref{fig:status2}).

\begin{figure}[ht!]
\centering
\begin{minipage}[c]{0.3\textwidth}
    \centering
    \includegraphics[width=\textwidth]{./qa4sm_graphics/bulk_barplot_status_0-\detokenize{$<ContentVars['NetcdfMetaVars']['val_dc_dataset0']>$}_and_1-\detokenize{$<ContentVars['NetcdfMetaVars']['val_dc_dataset1']>$}.png}
    \subcaption{} % (a)
\end{minipage}
\begin{minipage}[c]{0.6\textwidth}
    \centering
    \includegraphics[width=\textwidth]{./qa4sm_graphics/bulk_overview_0-\detokenize{$<ContentVars['NetcdfMetaVars']['val_dc_dataset0']>$}_and_1-\detokenize{$<ContentVars['NetcdfMetaVars']['val_dc_dataset1']>$}_status.png}
    \subcaption{} % (b)
\end{minipage}
\caption{Number of successful test cases and errors raised (a), and their respective location (b) in the validation run against \detokenize{$<ContentVars['ConfigVars']['DS0']['name']>$}}
\label{fig:status2}
\end{figure}

\noindent Validation metrics are computed from common observation time stamps between \ac{SMOS} \ac{L2} and ERA5-Land.
A \detokenize{$<ContentVars['ConfigVars']['temporal_matching']>$}-hour temporal matching window is applied to each \detokenize{$<ContentVars['ConfigVars']['temporal_reference']>$} observation, to find the closest data point in the other dataset.
The resulting sample size after matching, which applies to all computed metrics, is shown in Fig. \ref{fig:nobs2}.


\begin{figure}[ht!]
\centering
\begin{minipage}[c]{0.3\textwidth}
    \centering
    \includegraphics[width=\textwidth]{./qa4sm_graphics/bulk_boxplot_n_obs.png}
    \subcaption{} % (a)
\end{minipage}
\begin{minipage}[c]{0.6\textwidth}
    \centering
    \includegraphics[width=\textwidth]{./qa4sm_graphics/bulk_overview_n_obs.png}
    \subcaption{} % (b)
\end{minipage}
\caption{Validation sample size distribution (a), and by validation location (b) in the validation run against ERA5-Land.}
\label{fig:nobs2}
\end{figure}

\noindent Fig. \ref{fig:Rbox2} shows the correlation coefficients (Pearson's R) (a, b) and p-values (c) between \ac{SMOS} \ac{L2} and ERA5-Land at all successfully tested locations.

\begin{figure}[ht!]
\centering
\begin{minipage}[c]{0.45\textwidth}
    \centering
    \includegraphics[width=\textwidth]{./qa4sm_graphics/bulk_boxplot_R.png}
    \subcaption{} % (a)
\end{minipage}%
\hfill
\begin{minipage}[c]{0.5\textwidth}
    \centering
    \includegraphics[width=\textwidth]{./qa4sm_graphics/bulk_overview_0-\detokenize{$<ContentVars['NetcdfMetaVars']['val_dc_dataset0']>$}_and_1-\detokenize{$<ContentVars['NetcdfMetaVars']['val_dc_dataset1']>$}_R.png}
    \subcaption{} % (b)
    \medskip
    \includegraphics[width=\textwidth]{./qa4sm_graphics/bulk_overview_0-\detokenize{$<ContentVars['NetcdfMetaVars']['val_dc_dataset0']>$}_and_1-\detokenize{$<ContentVars['NetcdfMetaVars']['val_dc_dataset1']>$}_p_R.png}
    \subcaption{} % (c)
\end{minipage}
\caption{Distribution of Pearson's R values across all locations (a), spatial overview of Pearson's R (b) and corresponding p-values (c) between common observations from \ac{SMOS} \ac{L2} and ERA5-Land}
\label{fig:Rbox2}
\end{figure}


